% GENERATED FILE. Edit canonical lessons, then run npm run generate:books.
% canonical-source-hash: fnv1a64:0364b864fa6f2ec2
% canonical-lessons: SA-C343-vadah, SA-C343-tarkah, SA-C343-vakta, SA-C343-srota, SA-C343-pratipadayati

\chapter{A Debate}
\label{ch:sa-a-debate}
\hlchaptermodality{343}

\begin{chapteropening}
\textbf{By the end of this chapter:} \emph{I can say who spoke in a debate and who listened.}

Complete the last lesson of chapter 343: I can say who spoke in a debate and who listened.
\end{chapteropening}

\section[\texorpdfstring{\sk{वादः} (vādaḥ)}{वादः (vādaḥ)}]{\sk{वादः} (vādaḥ) — a debate; a thesis}

\label{lesson:SA-C343-vadah}



\begin{quote}
Before the new one: say the Sanskrit for as a result, then the Sanskrit for argues for.
\end{quote}

\subsection*{You'll want to know: \sk{वादः}}
\textbf{\sk{वादः}} — \emph{vādaḥ} — \textquotedblleft{}a debate; a thesis\textquotedblright{}.

\begin{quote}
\sk{श्वः} \sk{सभायां} \sk{वादः} \sk{भविष्यति।}
\end{quote}

\begin{quote}
\emph{śvaḥ sabhāyāṃ vādaḥ bhaviṣyati.}
\end{quote}

\begin{quote}
There will be a debate in the assembly tomorrow.
\end{quote}

A masculine noun.

\begin{cousinweb}
\textbf{\sk{वादः}} (\emph{vādaḥ}) comes from the root √vad, \textquotedblleft{}to speak\textquotedblright{}; MW gives \emph{vāda} the senses \textquotedblleft{}a thesis, an argument\textquotedblright{} and \textquotedblleft{}discussion, controversy, dispute\textquotedblright{} (MW).
\end{cousinweb}

\begin{grammarlens}[title={What you've built}]
One more word for giving reasons and opinions.
\end{grammarlens}

\subsection*{Your turn}
\begin{itemize}
\raggedright
  \item \emph{Say it:} \emph{vādaḥ}
  \item \emph{Say it:} \emph{vādaḥ}, once more
  \item \emph{Say it:} \emph{śvaḥ sabhāyāṃ vādaḥ bhaviṣyati.} — There will be a debate in the assembly tomorrow.
  \item \emph{Recall:} say the Sanskrit for as a result, then the Sanskrit for argues for, then say \emph{vādaḥ} again
\end{itemize}

\begin{tcolorbox}[breakable,colback=teal!4,colframe=teal!35!black,title={Before you move on}]
What does \sk{वादः} mean? (\textquotedblleft{}A debate; a thesis\textquotedblright{}.) Say it once more.
\end{tcolorbox}
\section[\texorpdfstring{\sk{तर्कः} (tarkaḥ)}{तर्कः (tarkaḥ)}]{\sk{तर्कः} (tarkaḥ) — reasoning}

\label{lesson:SA-C343-tarkah}



\begin{quote}
Before the new one: say the Sanskrit for argues for, then the Sanskrit for a debate.
\end{quote}

\subsection*{You'll want to know: \sk{तर्कः}}
\textbf{\sk{तर्कः}} — \emph{tarkaḥ} — \textquotedblleft{}reasoning\textquotedblright{}.

\begin{quote}
\sk{वादे} \sk{तर्कः} \sk{आवश्यकः।}
\end{quote}

\begin{quote}
\emph{vāde tarkaḥ āvaśyakaḥ.}
\end{quote}

\begin{quote}
In a debate, reasoning is necessary.
\end{quote}

A masculine noun.

\begin{cousinweb}
\textbf{\sk{तर्कः}} (\emph{tarkaḥ}) comes from the root √tark, \textquotedblleft{}to conjecture, to reason\textquotedblright{}; MW gives \emph{tarka} the senses \textquotedblleft{}conjecture\textquotedblright{} and \textquotedblleft{}reasoning, inquiry\textquotedblright{} (MW).
\end{cousinweb}

\begin{grammarlens}[title={What you've built}]
One more word for giving reasons and opinions.
\end{grammarlens}

\subsection*{Your turn}
\begin{itemize}
\raggedright
  \item \emph{Say it:} \emph{tarkaḥ}
  \item \emph{Say it:} \emph{tarkaḥ}, once more
  \item \emph{Say it:} \emph{vāde tarkaḥ āvaśyakaḥ.} — In a debate, reasoning is necessary.
  \item \emph{Recall:} say the Sanskrit for argues for, then the Sanskrit for a debate, then say \emph{tarkaḥ} again
\end{itemize}

\begin{tcolorbox}[breakable,colback=teal!4,colframe=teal!35!black,title={Before you move on}]
What does \sk{तर्कः} mean? (\textquotedblleft{}Reasoning\textquotedblright{}.) Say it once more.
\end{tcolorbox}
\section[\texorpdfstring{\sk{वक्ता} (vaktā)}{वक्ता (vaktā)}]{\sk{वक्ता} (vaktā) — a speaker}

\label{lesson:SA-C343-vakta}



\begin{quote}
Before the new one: say the Sanskrit for a debate, then the Sanskrit for reasoning.
\end{quote}

\subsection*{You'll want to know: \sk{वक्ता}}
\textbf{\sk{वक्ता}} — \emph{vaktā} — \textquotedblleft{}a speaker\textquotedblright{}.

\begin{quote}
\sk{वक्ता} \sk{स्पष्टं} \sk{वदेत्।}
\end{quote}

\begin{quote}
\emph{vaktā spaṣṭaṃ vadet.}
\end{quote}

\begin{quote}
A speaker should speak clearly.
\end{quote}

A masculine noun.

\begin{cousinweb}
\textbf{\sk{वक्ता}} (\emph{vaktā}) is built on the root √vac, \textquotedblleft{}to speak\textquotedblright{}, with the agent ending \emph{-tṛ}; MW gives \emph{vaktṛ} \textquotedblleft{}one who speaks, a speaker\textquotedblright{} (MW).
\end{cousinweb}

\begin{grammarlens}[title={What you've built}]
One more word for giving reasons and opinions.
\end{grammarlens}

\subsection*{Your turn}
\begin{itemize}
\raggedright
  \item \emph{Say it:} \emph{vaktā}
  \item \emph{Say it:} \emph{vaktā}, once more
  \item \emph{Say it:} \emph{vaktā spaṣṭaṃ vadet.} — A speaker should speak clearly.
  \item \emph{Recall:} say the Sanskrit for a debate, then the Sanskrit for reasoning, then say \emph{vaktā} again
\end{itemize}

\begin{tcolorbox}[breakable,colback=teal!4,colframe=teal!35!black,title={Before you move on}]
What does \sk{वक्ता} mean? (\textquotedblleft{}A speaker\textquotedblright{}.) Say it once more.
\end{tcolorbox}
\section[\texorpdfstring{\sk{श्रोता} (śrotā)}{श्रोता (śrotā)}]{\sk{श्रोता} (śrotā) — a listener}

\label{lesson:SA-C343-srota}



\begin{quote}
Before the new one: say the Sanskrit for reasoning, then the Sanskrit for a speaker.
\end{quote}

\subsection*{You'll want to know: \sk{श्रोता}}
\textbf{\sk{श्रोता}} — \emph{śrotā} — \textquotedblleft{}a listener\textquotedblright{}.

\begin{quote}
\sk{श्रोता} \sk{शान्तः} \sk{तिष्ठेत्।}
\end{quote}

\begin{quote}
\emph{śrotā śāntaḥ tiṣṭhet.}
\end{quote}

\begin{quote}
A listener should stay quiet.
\end{quote}

A masculine noun.

\begin{cousinweb}
\textbf{\sk{श्रोता}} (\emph{śrotā}) is built on the root √śru, \textquotedblleft{}to hear\textquotedblright{}, with the agent ending \emph{-tṛ}; MW gives \emph{śrotṛ} \textquotedblleft{}one who hears, a hearer\textquotedblright{} (MW).
\end{cousinweb}

\begin{grammarlens}[title={What you've built}]
One more word for giving reasons and opinions.
\end{grammarlens}

\subsection*{Your turn}
\begin{itemize}
\raggedright
  \item \emph{Say it:} \emph{śrotā}
  \item \emph{Say it:} \emph{śrotā}, once more
  \item \emph{Say it:} \emph{śrotā śāntaḥ tiṣṭhet.} — A listener should stay quiet.
  \item \emph{Recall:} say the Sanskrit for reasoning, then the Sanskrit for a speaker, then say \emph{śrotā} again
\end{itemize}

\begin{tcolorbox}[breakable,colback=teal!4,colframe=teal!35!black,title={Before you move on}]
What does \sk{श्रोता} mean? (\textquotedblleft{}A listener\textquotedblright{}.) Say it once more.
\end{tcolorbox}
\section[\texorpdfstring{\sk{प्रतिपादयति} (pratipādayati)}{प्रतिपादयति (pratipādayati)}]{\sk{प्रतिपादयति} (pratipādayati) — sets out (a case); explains}

\label{lesson:SA-C343-pratipadayati}



\begin{quote}
Before the new one: say the Sanskrit for a speaker, then the Sanskrit for a listener.
\end{quote}

\subsection*{You'll want to know: \sk{प्रतिपादयति}}
\textbf{\sk{प्रतिपादयति}} — \emph{pratipādayati} — \textquotedblleft{}sets out (a case); explains\textquotedblright{}.

\begin{quote}
\sk{अरुणः} \sk{तर्केण} \sk{मतं} \sk{प्रतिपादयति।}
\end{quote}

\begin{quote}
\emph{aruṇaḥ tarkeṇa mataṃ pratipādayati.}
\end{quote}

\begin{quote}
Arun sets out his view with reasoning.
\end{quote}

Past (he …): \sk{सः} \sk{मतं} \sk{प्रतिपादितवान्}, \emph{saḥ mataṃ pratipāditavān}.

\begin{cousinweb}
\textbf{\sk{प्रतिपादयति}} (\emph{pratipādayati}) is built on \emph{prati}, \textquotedblleft{}towards\textquotedblright{}, and the root √pad, \textquotedblleft{}to go\textquotedblright{}, in its causative form; MW gives the senses \textquotedblleft{}to establish, to prove, to set forth, to explain\textquotedblright{} (MW).
\end{cousinweb}

\begin{grammarlens}[title={What you've built}]
One more word for giving reasons and opinions.
\end{grammarlens}

\subsection*{Your turn}
\begin{itemize}
\raggedright
  \item \emph{Say it:} \emph{pratipādayati}
  \item \emph{Say it:} \emph{pratipādayati}, once more
  \item \emph{Say it:} \emph{aruṇaḥ tarkeṇa mataṃ pratipādayati.} — Arun sets out his view with reasoning.
  \item \emph{Recall:} say the Sanskrit for a speaker, then the Sanskrit for a listener, then say \emph{pratipādayati} again
\end{itemize}

\begin{tcolorbox}[breakable,colback=teal!4,colframe=teal!35!black,title={Before you move on}]
What does \sk{प्रतिपादयति} mean? (\textquotedblleft{}Sets out (a case); explains\textquotedblright{}.) Say it once more.
\end{tcolorbox}
